\section{The represented, missing, and full physical maps}
\label{sec:tpc60-physical-interface}

We first place all three outputs on one coefficient space.  This prevents
two common changes of problem: choosing the missing vector independently
of the represented vector, and comparing norms formed from different
atomic normalizations.

\subsection{Fixed carriers and the literal high partition}

Let \(\cW_0\) be the complete declared Walsh slice of TPC--59, after
removing only coordinates which are identically inactive throughout the
fixed bounded coefficient class.  Throughout, \(y\) is TPC--59
very-high admissible:
\begin{equation}
 y>\max\{|h_0|,D_+,J_+,\Delta_M,\sqrt{N_+}\},
 \qquad
 \frac{D_+N_+}{y}<D_-N_-,
 \label{eq:tpc60-very-high-admissibility}
\end{equation}
and every retained source row has its inherited unique representation
\(m=d_m\ell_m\) in the declared one-side row system
\citep{WangTPC59}.  Let \(\cW_H(y)\) be the resulting very-high part.
The fixed-shift terminal census supplies an injection
\begin{equation}
 \iota:\cU_{h_0,>y}^{\rm term}\longrightarrow\cW_H(y).
 \label{eq:tpc60-terminal-Walsh-injection}
\end{equation}
Set
\begin{equation}
 \cW_B:=\cW_0\setminus\cW_H(y),
 \qquad
 \cW_R:=\Ran\iota,
 \qquad
 \cW_M:=\cW_H(y)\setminus\cW_R.
 \label{eq:tpc60-BRM-carriers}
\end{equation}
The letters \(B,R,M\) stand for base, represented high, and missing high.
They form the disjoint partition
\begin{equation}
 \cW_0=\cW_B\,\dot\cup\,\cW_R\,\dot\cup\,\cW_M.
 \label{eq:tpc60-W0-partition}
\end{equation}

Let \(\cK_X\) be the raw finite row-coefficient Hilbert space.  For each
Walsh atom \(w\), literal factorization gives
\begin{equation}
 b_w(\zeta)=\beta_w\zeta_{m(w)},
 \qquad \zeta\in\cK_X,
 \label{eq:tpc60-literal-atom-linear}
\end{equation}
with \(\beta_w\) fixed by the declared row, orbit, cutoff, and profile
data.  No favorable atom-dependent coefficient is introduced.  The raw
atomic diagonals are
\begin{equation}
 \cD_E(\zeta):=\sum_{w\in\cW_E}|b_w(\zeta)|^2,
 \qquad E\in\{B,R,M\}.
 \label{eq:tpc60-BRM-diagonals}
\end{equation}
In the abstract atomic notation used below, the missing lift is normalized
by the same physical amplitudes, so
\begin{equation}
 \cD_{\rm miss}(\zeta)=\cD_M(\zeta).
 \label{eq:tpc60-missing-diagonal-identification}
\end{equation}
Therefore
\begin{equation}
 \boxed{
 \cD_0=\cD_B+\cD_R+\cD_M,}
 \qquad
 \cD_R=\sum_k\cD_k
 \label{eq:tpc60-exact-raw-partition}
\end{equation}
on the critical high face, with the dyadic cells of TPC--59.  This is a
Pythagorean identity in the atomic space; it is not an additivity statement
for the coherently grouped output.

\subsection{Recovery of the physical labels}

For a nonzero high Walsh atom \(w\), the literal fixed-shift dictionary
recovers a unique triple \((m,j,\gamma)\), the target
\begin{equation}
 n=mj+h_0,
 \label{eq:tpc60-target-n}
\end{equation}
and, because \(y>\sqrt{N_+}\), the unique factorization
\begin{equation}
 p=P^+(n)>y,
 \qquad r=\frac np,
 \qquad s=d_mr.
 \label{eq:tpc60-prs-recovery}
\end{equation}
The native output key is denoted \(i=i(w)\); on the displayed left slice
it includes the side, orbit tag, and row-time data, and the right slice is
handled by orthogonal direct sum.  We write
\begin{equation}
 q(w)=(i(w),s(w)).
 \label{eq:tpc60-residual-key}
\end{equation}
In particular, the native key \(i\) fixes the side, the opposite row, and
the orbit time \(j\).  All of these labels are recovered from the physical
atom.  They are not
new choices made after seeing \(\zeta\).

For \(E\in\{R,M\}\), define the sign-compatible high coefficient by
\begin{equation}
 z_E(i,s)
 :=-\sum_{\substack{w\in\cW_E\\q(w)=(i,s)}}b_w.
 \label{eq:tpc60-zRM}
\end{equation}
Let
\begin{equation}
 (C_\mu z)_i:=\sum_s\mu(s)z(i,s),
 \qquad
 h_R:=C_\mu z_R,
 \qquad h_M:=C_\mu z_M.
 \label{eq:tpc60-hRM}
\end{equation}
The minus sign in \eqref{eq:tpc60-zRM} is the fixed high-character sign;
the remaining \(\mu(s)\) is the low-label gauge.  Thus these formulas use
the TPC--58 sign convention and do not apply a second M\"obius sign.

For the base carrier write
\begin{equation}
 (v_0)_i
 :=\sum_{\substack{w\in\cW_B\\i(w)=i}}
       \mu(u_w)b_w,
 \label{eq:tpc60-base-vector}
\end{equation}
where \(u_w\) is its low Walsh label.  The precise inherited support masks
are already contained in \(\cW_B\).

\begin{proposition}[Literal three-map identity]
\label{prop:tpc60-literal-three-map}
There are linear maps
\begin{equation}
 A_X,M_X,B_X:\cK_X\longrightarrow\cH_{\rm out}
 \label{eq:tpc60-three-map-domains}
\end{equation}
such that, for every physical row coefficient vector \(\zeta\),
\begin{align}
 A_X\zeta&=v_0(\zeta)+h_R(\zeta),
 \label{eq:tpc60-represented-map}\\
 M_X\zeta&=h_M(\zeta),
 \label{eq:tpc60-missing-map}\\
 B_X\zeta&=v_0(\zeta)+h_R(\zeta)+h_M(\zeta)
            =(A_X+M_X)\zeta.
 \label{eq:tpc60-full-map}
\end{align}
Moreover \(h_R=C_{\mu,H}T_H^{\rm rep}\zeta\) is exactly the TPC--58 high
map restricted to \(\cW_R=\Ran\iota\), while \(h_M\) is the same physical
grouping restricted to \(\cW_M\).  Consequently
\begin{equation}
 \cN_R(\zeta):=\|A_X\zeta\|_2^2,
 \qquad
 \cN_{\rm full}(\zeta):=\|B_X\zeta\|_2^2
 \label{eq:tpc60-represented-full-energies}
\end{equation}
are the represented and full native affine energies on one and the same
declared slice.
\end{proposition}

\begin{proof}
Linearity follows from \eqref{eq:tpc60-literal-atom-linear} and the fixed
carrier sums.  The gauge alignment of TPC--58 identifies the represented
terminal output with \(C_{\mu,H}T_H^{\rm rep}\zeta\).  Restricting the
identical atom-to-native map to the complementary high set defines \(M_X\).
The disjoint coordinate partition \eqref{eq:tpc60-W0-partition} then gives
the last identity before any output norm is taken.
\end{proof}

\begin{remark}[Why the atom partition does not imply output monotonicity]
The equality \(\cD_0=\cD_B+\cD_R+\cD_M\) is taken before native grouping.
After grouping,
\[
 \cN_{\rm full}-\cN_R
 =\|h_M\|^2+2\operatorname{Re}\langle A_X\zeta,h_M\rangle.
\]
The cross term has no fixed sign.  Thus restoring a disjoint atomic
coordinate set can decrease the physical output norm.
\end{remark}

\subsection{Prime-resolved alias structure}

To locate the possible cancellation more precisely, decompose
\begin{equation}
 h_R=-\sum_{p>y}v_p^R,
 \qquad
 h_M=-\sum_{p>y}v_p^M,
 \label{eq:tpc60-prime-decompositions}
\end{equation}
where, for \(E\in\{R,M\}\),
\begin{equation}
 (v_p^E)_i
 :=\sum_{\substack{w\in\cW_E\\i(w)=i,\ p(w)=p}}
       \mu(s(w))b_w.
 \label{eq:tpc60-prime-vector-components}
\end{equation}
On the active very-high carrier, squarefree support gives
\(|\mu(s(w))|=1\).  Hence no modulus is lost when the raw atom is placed in
a prime-resolved component.

\begin{proposition}[Same-prime separation and unequal-prime alias]
\label{prop:tpc60-same-prime-separation}
For every fixed pair \((i,p)\), there is at most one active high Walsh atom.
Consequently
\begin{equation}
 \langle v_p^R,v_p^M\rangle=0
 \label{eq:tpc60-same-prime-orthogonality}
\end{equation}
for every \(p\), and
\begin{equation}
 \sum_p\bigl(\|v_p^R\|^2+\|v_p^M\|^2\bigr)
 =\cD_R+\cD_M.
 \label{eq:tpc60-prime-resolved-Pythagoras}
\end{equation}
The exact reassembly defect is
\begin{align}
 \cN_{\rm full}-\cN_R
 ={}&\left\|\sum_qv_q^M\right\|^2
 -2\operatorname{Re}\sum_q\langle v_0,v_q^M\rangle
 \notag\\
 &+2\operatorname{Re}
   \sum_{\substack{p,q>y\\p\ne q}}
       \langle v_p^R,v_q^M\rangle.
 \label{eq:tpc60-exact-prime-alias-defect}
\end{align}
Thus represented--missing cancellation cannot occur at the same prime;
the remaining dangers are base--missing interaction and unequal-prime
alias at a common native output.
\end{proposition}

\begin{proof}
Fix \((i,p)\).  Since \(i\) fixes \(j\), suppose that \(p\) meets two
source rows \(m,m'\).  Then \(p\mid(m-m')j\).  The alternative \(p\mid j\)
is excluded by \(p>J_+\); hence \(p\mid m-m'\).  But
\(|m-m'|\le\Delta_M<p\), so \(m=m'\), and the inherited unique row
representation recovers the same full source atom.  This is exactly the
TPC--59 critical-cutoff fixed-\((j,p)\) argument
\citep{WangTPC59}.  Thus a fixed \((i,p)\) contains at most one high atom.
Such an atom belongs to exactly one of
\(\cW_R,\cW_M\), proving disjoint coordinate support and
\eqref{eq:tpc60-same-prime-orthogonality}.  Summing the coordinate squares
gives \eqref{eq:tpc60-prime-resolved-Pythagoras}.  Finally expand
\[
 \|v_0-\sum_pv_p^R-\sum_qv_q^M\|^2
 -\|v_0-\sum_pv_p^R\|^2
\]
and delete the same-prime cross terms using
\eqref{eq:tpc60-same-prime-orthogonality}.
\end{proof}

The proposition is useful but does not close reassembly.  Different high
primes can land at the same native output, and degree two already permits
exact cancellation.  Prime resolution is therefore a diagnostic layer,
not a substitute for the full native Gram.

\subsection{The six-item physical audit}

The construction above passes the following interface audit.
\begin{enumerate}[leftmargin=2.2em]
 \item The coefficients are the literal \(b_w=\beta_w\zeta_{m(w)}\).
 \item One fixed physical \(h_0\) appears in \(n=mj+h_0\).
 \item Every atomic norm is the raw counting norm; no fiber average is
       inserted.
 \item The atom representation is nonduplicated, and the row space removes
       only identically null columns.  A quotient by \(\Ker A_X\) is not
       taken unless the full map actually descends to it.
 \item The active supports are exactly \(\cW_B,\cW_R,\cW_M\) inside the
       declared Walsh carrier; coefficient-dependent zeros do not redefine
       them.
 \item Any eventual reassembly loss remains an explicit term in the strict
       endpoint budget \(<1/400\).
\end{enumerate}
